\section{Discussion}

Marine Race Arena is best understood as evaluation infrastructure rather than a controller paper. This distinction is important. A new controller can be impressive but difficult to compare if every paper uses different scenarios, metrics, and hidden assumptions. A benchmark arena has value when it allows many controllers to be tested under the same rules. The proposed design therefore emphasizes explicit configurations, official sensor restrictions, repeatable tasks, and referee-controlled scoring.

The main limitation is that the current paper is simulation-only. This is not a minor limitation: underwater sim-to-real transfer is hard because visual appearance, hydrodynamics, tether effects, actuation limits, acoustic propagation, lighting, and sensor noise can differ substantially between simulation and water-tank experiments. For this reason, the current contribution should not be framed as a validated predictor of real underwater racing performance. It should be framed as a reproducible development and evaluation layer that can later be connected to physical tests.

A second limitation concerns gate validation. The current center-point rule is deterministic and simple, but it is not a full vehicle-body clearance model. A vehicle can have a valid center crossing while still passing too close to a physical bar. Clearance margins partially address this, but future versions should support footprint- or mesh-aware validation.

A third limitation is controller maturity. The acoustic baseline and rule-based vision baseline are useful starting points, but the final paper should include systematic runs before claiming robustness. In particular, currents are the decisive underwater contribution: if the arena includes currents but the baselines are reported only on clean tracks, the evaluation will look incomplete. The strongest final version of the paper should therefore include current-gate results, even if the baseline fails often. Failure under currents is still valuable if it exposes why the benchmark is needed.

Future work should extend the arena in four directions: richer hydrodynamic profiles, sensor-degradation models, full-body gate validation, and standardized submission containers for external controllers. A later sim-to-real extension could use a small tank and a BlueROV2-style platform to validate whether the same referee and controller interface can support real-water trials.
